\phantomsection
\part{用語の定義}

\begin{abstract}
  この章では，数学の学習を始める前に，定義・命題・定理・補題・系などの基本的な用語について理解を深める．「『命題』という言葉は知っているけど，『定理』とはどう違うの？」という疑問を持ったり，「『定義』と『定理』の違いがあやふやなままで，よくわからない」と感じる学生は多い．まず，日常的な例と数学的な例をもとに「定義」の性質を確認する．次に，命題・定理・補題・系という用語を整理し，これらのラベルの使い分けについて注意を述べる．最後に，「かつ・または」「とる・おく・固定する」「十分大きい」など，数学に特有の言い回しを具体例とともに確認しよう．
\end{abstract}

\phantomsection
\section{定義・命題・定理・補題・系}

\phantomsection
\subsection{定義}

数学用語には細かな意味がある．概要にも述べたように，「定義」と「定理」の違いなどをはっきりと認識しておくことは，
数学の学習を進める上で堅固たる「理解の基盤」となる．まずは「定義」の意味の確認からはじめよう：

\emph{定義}\index{ていぎ@定義}（\emph{definition}）とは，用語の意味を明確に述べたものであり，\textbf{Def}と略記される．
要するに「定義」とは，「この数学的概念（用語）はこのような意味である」という取り決めである．
数学のお話にうつる前に，日常的な「定義」の例を見てみよう．

\begin{example}{アレニウスによる酸・塩基の定義}{アレニウスによる酸・塩基の定義}
  水溶液中で水素イオン\ce{H+}を放出する物質を酸，\ce{OH-}を放出する物質を塩基とする．
\end{example}

\begin{example}{「民主主義」の定義}{「民主主義」の定義}
  民主主義とは，国民が政治の意思決定に直接または間接的に参加する政治体制のことをいう．
\end{example}

どちらもわかりやすい例であり，「すでに知っていること」と思われる読者もいるかもしれない．
しかし，議論を行ううえで，定義が曖昧なままだと不都合が生じることがある．

\begin{example}{「速度」と「速さ」の違い}{「速度」と「速さ」の違い}
直線上を運動する物体の位置を$x$とすると，その速度$v$\footnote{空間内の運動では，速度を$\vec{v}$や$\bm{v}$などと表すこともある．一方，大学の数学や物理では，文脈から明らかな場合にはベクトルとスカラーを表記上区別しないことも多い．}は
\[
v=\frac{dx}{dt}
\]
と定義される．ここで，「速度」と「速さ」を同じ意味で用いたとしよう．

速度$v$には正負があり，その符号は物体がどちらの向きに運動しているかを表す．一方，速さは運動の方向を考えず，速度の大きさ
\[
\abs{v}
\]
で表される．

たとえば，速度が$v=-3$である物体の速さは$3$である．ここで速度と速さを混同して$v=3$としてしまうと，物体が実際とは反対向きに運動していることになってしまう．
このように，用語の意味を曖昧にすると必要な情報が失われ，まったく異なる結論に至ることがある．
\end{example}

\begin{example}{「平均」と「期待値」の違い}{「平均」と「期待値」の違い}
  ここで比較するのは，\emph{算術平均}（標本平均）と期待値である．算術平均とは，実際に得られたデータをならした値である．
  たとえば，$n$個のデータ $x_1,x_2,\dots,x_n$ の算術平均は
  \[
    \frac{x_1+x_2+\cdots+x_n}{n}
  \]
  で与えられる．

  算術平均の定義に基づくと，$1$，$4$，$10$ という3個のデータの算術平均は
  \[
    \frac{1+4+10}{3}=5
  \]
  である．

  一方，期待値とは，各値にその値が出る確率をかけて足し合わせたものである．
  たとえば，ある確率変数 $X$ が $1$，$4$，$10$ をそれぞれ
  $\frac{1}{6}$，$\frac{1}{6}$，$\frac{2}{3}$ の確率でとるとき，
  その期待値は
  \[
    E[X]=1\cdot\frac{1}{6}+4\cdot\frac{1}{6}+10\cdot\frac{2}{3}=7.5
  \]
  である．

  このように，算術平均は実際のデータから求める値であり，
  期待値は確率を用いて求める値である\footnote{ただし，確率論では期待値のことを「平均」とよぶこともある．「平均」という語は文脈により指すものが異なるので，どの意味で用いられているかに注意したい．}．
\end{example}

\exref{「速度」と「速さ」の違い}や\exref{「平均」と「期待値」の違い}からもわかるように，
数学の学習において「定義」をあやふやにしてはならない．

しかし，同じ事柄について，2つ以上の定義の形式があることもある\footnote{「書籍によって採用している定義が違う」というのもよくあることである．}．次は数学における例を見てみよう：

\phantomsection
\begin{example}{「絶対値」の定義}{「絶対値」の定義}
  $ x \in \mathbb{R}$に対して，$x$の絶対値を$\abs{x}$とかき，次のように定義する：
  \[
    \abs{x} \coloneqq \max \{ x , -x \}
  \]
  この定義は以下のような形式で表現してもよい：
  \[
    \abs{x} \coloneqq  \sqrt{x^2}
  \]
\end{example}

つまり，$x\in\mathbb{R}$の絶対値$\abs{x}$は，$\abs{x} \coloneqq \max \{ x, -x \}$としてもよい一方，
$\abs{x} \coloneqq \sqrt{x^2}$と定義して問題ないのである．どちらの定義を採用してもよく，もちろんこの他の定義を採用してもよい\footnote{$x\geqq 0$のとき$\abs{x} \coloneqq x$，$x<0$のとき$\abs{x} \coloneqq -x$とする方法などがある．}

また，\exref{「絶対値」の定義}のように，同じ概念にふたつの定義の形式があるとき，
「どちらの形式を採用しても同じ概念が定まる」こと，すなわち「ふたつの定義が同値である」ということは，
証明を要する主張である\footnote{「証明」とはなにか，どのように書くのかについては\partref{証明}で詳しく扱う．
本章では「定義や既知の事実から，筋道を立てて主張が正しいことを確かめる手続き」という程度の理解でよく，
本章に現れる証明は書き方の雰囲気をつかむための例と思って読み進めてほしい．細部は\partref{証明}を読んだあとで改めて確認すればよい．}．
片方を定義として採用した場合には，もう片方は「証明すべき命題」という扱いになるのである．
実際に確かめてみよう．
 
\begin{prop}{絶対値のふたつの定義の同値性}{絶対値のふたつの定義の同値性}
  任意の$x \in \mathbb{R}$に対して，以下の等式が成り立つ：
  \[
    \max \{ x , -x \} = \sqrt{x^2}
  \]
  ただし，$a \geqq 0$に対して$\sqrt{a}$は「$2$乗すると$a$になる非負の実数」を表す\footnote{この記法が正当であるためには，「$2$乗すると$a$になる非負の実数」がただひとつ存在することが必要である．このことについては，のちの「一意性」の節でも触れる．}．
\end{prop}
 
\begin{leftbar}
  \begin{proof}
    $x$の符号で場合分けする．
      \proofstep{$x \geqq 0$の場合}
            $-x \leqq 0 \leqq x$により，$\max \{ x,-x \} = x$である．
            一方，$x \geqq 0$かつ$x^2 = x \cdot x$なので，$\sqrt{x^2}=x$である．
            よって両辺は一致する．
      \proofstep{$x < 0$の場合}
            $x < 0 < -x$により，$\max \{ x,-x \} = -x$である．
            一方，$-x > 0$かつ$x^2 = (-x)^2$なので，$\sqrt{x^2}=-x$である．
            よって両辺は一致する．
            
            \smallskip
    以上により，任意の$x \in \mathbb{R}$に対して主張の等式が成り立つ．これが証明すべきことであった．
  \end{proof}
\end{leftbar}

このように，\prref{絶対値のふたつの定義の同値性}に証明を与えることで，$x\in \mathbb{R}$の絶対値$\abs{x}$の定義として
\begin{enumerate*}[label=(\arabic*)]
  \item $\abs{x} \coloneqq \max \{x,-x\}$，
  \item $\abs{x} \coloneqq \sqrt{x^2}$
\end{enumerate*}
のどちらを採用してもよいことがわかった．


次に，具体例を通して「定義の性質」についてみていこう．具体的には「定義を設ける上での慣習」について確認をする．
線型代数の講義で学ぶことであるが，逆行列の定義は以下のようになる：

\begin{example}{「逆行列」の定義}{「逆行列」の定義}
  正方行列$A$に対して，$BA = AB =E$となるような$B$が存在するとき，このような$B$を$A$の逆行列という．
\end{example}


ここで注意するのは \emph{数学では定義は最小限の情報にとどめることが慣習となっている}ということである．たとえば，\exref{「逆行列」の定義}から，以下の命題\footnote{「命題」については，のちほど詳しくみていくとする．ここでは「証明を与えるべきである主張」という理解でよい．}が成り立ち，その証明は容易である．

\phantomsection
\begin{prop}{逆行列の一意性}{逆行列の一意性}
  正方行列$A$に対して，$A$の逆行列は存在するとすればただひとつである．
\end{prop}

\begin{tleftbar}
  \begin{proof}
    $A$の逆行列が$B$，$C$であるとすると，
    \[
      AB=BA=E,\quad AC=CA=E
    \]
    が成り立つ．このとき，
    \[
      B=BE=B(AC)=(BA)C=EC=C
    \]
    が成り立つ．よって，$B=C$であり，逆行列の一意性が示された．これが証明すべきことであった．
  \end{proof}
\end{tleftbar}

このことから，\exref{「逆行列」の定義}でとりあげた逆行列の定義をもっと詳しく

\begin{quote}
  正方行列$A$に対して，$BA = AB =E$となるような$B$が存在するとき，このような$B$はただひとつで，$B$を$A$の逆行列という．
  $B$は一意に存在するので，これを$A^{-1}$とかく\footnote{一意に存在することがわかっていないと，$A^{-1}$のような記法で表すことはためらわれる．}．もちろん
  \[
    AA^{-1}=A^{-1}A=E
  \]
  が成り立つ．
\end{quote}
と，「逆行列の一意性を定義に含めてもいいのではないか．」と主張する人がいるかもしれない．
ただ，数学では「定義は最小限の情報にとどめ，そこから導かれる主張を命題として証明する」という慣習があり，
なにが「定義」で，なにが「証明すべきこと」であるかはっきりさせることが多い\footnote{ただ，実際はここで取り上げた逆行列の定義も情報過多である．線型代数で学ぶことになるが，$AB=E$と$BA=E$のどちらか片方の式のみで定義してよいからである．}．

もうひとつ，定義について注意しておきたいことがある．
それは，\emph{適切に与えられた定義は，通常，真偽を問う対象ではなく，理論全体をうまく構築できるように「設定」される}ということである\footnote{「通常」と述べたのは，「対象を定められていない不適切な定義」がありうるからである．このことは，のちに述べるwell-definedの概念と関わる．}．
 
\phantomsection
\begin{example}{「$1$は素数でない」という約束}{「1は素数でない」という約束}
  素数は「$1$とその数自身以外に正の約数を持たない，$2$以上の自然数」と定義され，$1$は素数から除外される．
  これは「$1$を素数に含めるとなにか矛盾が生じるから」ではない．素因数分解の一意性
  \begin{quote}
    $2$以上の任意の自然数は，順序の違いを除いてただひととおりに素数の積として表される．
  \end{quote}
  という定理をきれいな形で述べるためである．もし$1$を素数に含めると，
  \[
    6 = 2 \cdot 3 = 1 \cdot 2 \cdot 3 = 1 \cdot 1 \cdot 2 \cdot 3 = \cdots
  \]
  となり，「ただひととおり」という主張が崩れてしまう\footnote{正確には，一意性の主張に「$1$の個数の違いを除いて」といった但し書きを毎回つける必要が生じる．}．
\end{example}

このように，「$1$は素数か」という問いに対する答えは，証明によって得られるのではなく，
「どのような定義を採用すると理論をうまく構築できるか」という観点から選ばれている．
$0\kaijou = 1$や$a^0 = 1$（$a \ne 0$）という約束も，既存の法則（$n\kaijou = n \cdot (n-1)\kaijou$や指数法則$a^{m+n} = a^m a^n$など）を例外なく記述できるようにする，という発想による．

また，新しい概念や関数を定義するときには，その定義がwell-defined\index{welldefined@well-defined}であることを確認する必要がある．
たとえば，有理数全体の集合を$\mathbb{Q}$とし，関数
\[
  f \colon \mathbb{Q} \to \mathbb{Z}
\]
を
\[
  f\left(\frac{a}{b}\right)=a+b
\]
によって定義したとしよう．このとき，
\[
  f\left(\frac12\right)=1+2=3
\]
である一方，
\[
  f\left(\frac24\right)=2+4=6
\]
となる．ところが，
\[
  \frac12=\frac24
\]
であるから，これは$f$が同じ有理数に対して異なる値を与えていることになる．一般に，
\[
  \frac12=\frac24=\frac36=\cdots
\]
のように，有理数は複数の表し方をもつ．同じ有理数をどのように表しても関数の値が一致しなければ，その規則は「有理数を定義域とする関数」を定義していることにはならない．
このような事例を踏まえ，対象の表し方によらず値が一意に定まることを，well-defined（よく定義されている）という．たとえば，
\[
g \colon \mathbb{Q} \to \mathbb{Q}
\]
を
\[
g\!\left(\frac{a}{b}\right)=\frac{2a}{b}
\]
と定めると，$g$はwell-definedである．

well-definedについては，\ref{subsec:写像の合成とwell-defined}項で再度確認をするので，ここではこの程度の説明に留めよう．

\phantomsection
\subsection{命題}

数学でいう\emph{命題}\index{めいだい@命題}（\emph{proposition}）という語には，「広い意味」と「狭い意味」のふたつの用法がある．
\emph{広い意味}では，真であるか偽であるかが定まる文を命題という\footnote{たとえば，「$1+1=2$」は真の命題であり，「$1+1=3$」は偽の命題である．真であっても偽であっても，真偽が定まってさえいれば命題とよぶことに注意しよう．}．
一方，\emph{狭い意味}では，「定理」「補題」「系」などと同じように，正しい数学の主張につけるラベルのひとつとして「命題」という語を用いる\footnote{前項の\prref{絶対値のふたつの定義の同値性}や\prref{逆行列の一意性}に「命題」という見出しをつけたのは，この用法である．}．

\begin{figure}[h]
  \centering
  \begin{tikzpicture}[
    scale=0.90,
    transform shape,
    node distance=5mm and 10mm
  ]
    \node[flowbox] (sent) {言葉・式};
    \node[flowq, below=of sent] (q0) {主張を表しているか？};
    \node[flowbox, right=12mm of q0] (nosent)
      {命題ではない\\ \scriptsize（「$1+1$」「りんご」など）};

    \node[flowq, below=of q0] (q1) {真偽が定まるか？};
    \node[flowbox, right=12mm of q1] (notprop)
      {命題ではない\\ \scriptsize（「この部屋は暑い」など）};

    \node[flowbox, flowhl, below=of q1] (broad) {命題（広い意味）};
    \node[flowq, below=of broad] (q2) {真か偽か？};

    \node[flowbox, left=10mm of q2, yshift=-13mm] (false)
      {偽の命題\\ \scriptsize（「$1+1=46$」など）};
    \node[flowbox, right=10mm of q2, yshift=-13mm] (true)
      {真の命題\\ \scriptsize（「$1+1=2$」など）};

    \node[flowq, below=of true] (q3) {証明を与え，\\見出しをつけるか？};

    \node[flowbox, flowhl, below=8mm of q3, xshift=-14mm] (narrow)
      {命題（狭い意味）};
    \node[flowlab, right=3mm of narrow] (thm) {定理};
    \node[flowlab, right=3mm of thm]    (lem) {補題};
    \node[flowlab, right=3mm of lem]    (cor) {系};

    \node[
      draw=appleGray,
      dashed,
      rounded corners=3pt,
      inner sep=5pt,
      fit=(narrow)(thm)(lem)(cor),
      label={[flowtxt]below:
        正しいと分かっている主張につける見出し（ラベル）}
    ] (labels) {};

    \draw[flowarr] (sent) -- (q0);
    \draw[flowarr] (q0) -- node[flowtxt, above] {いいえ} (nosent);
    \draw[flowarr] (q0) -- node[flowtxt, right] {はい} (q1);

    \draw[flowarr] (q1) -- node[flowtxt, above] {いいえ} (notprop);
    \draw[flowarr] (q1) -- node[flowtxt, right] {はい} (broad);

    \draw[flowarr] (broad) -- (q2);
    \draw[flowarr] (q2) -|
      node[flowtxt, pos=0.25, above] {偽} (false);
    \draw[flowarr] (q2) -|
      node[flowtxt, pos=0.25, above] {真} (true);

    \draw[flowarr] (true) -- (q3);
    \draw[flowarr] (q3) --
      node[flowtxt, right] {はい} (labels.north -| q3);
  \end{tikzpicture}
  \caption{「命題」のふたつの意味}
  \label{fig:命題のふたつの意味}
\end{figure}

本項では「広い意味での命題」について解説する．そして，本書ではどちらの意味であるかが明らかな場合は，単に「命題」とかくと約束する．

「広い意味での命題」は「真の命題」と「偽の命題」のいずれかであり，このようにふたつの値を用いる論理体系を\emph{二値論理}\index{にちろんり@二値論理}とよぶ．
以下では，命題が真であることを$\mathrm{T}$，偽であることを$\mathrm{F}$と表す．

命題の定義にあてはまらないとされる主張には

\begin{enumerate}[(1)]
\item 用語や対象の定義が曖昧なもの
  \label{enu:定義が曖昧なもの}
\item 文の意味や解釈が曖昧なもの
  \label{enu:意味が曖昧なもの}
\item 真偽を問うことのできない文
  \label{enu:真偽を問えないもの}
\item 真偽が一意に定まらないもの
  \label{enu:真偽が定まらないもの}
\item パラドックス
  \label{enu:パラドックス}
\end{enumerate}

などがある．\ref{enu:定義が曖昧なもの}，\ref{enu:意味が曖昧なもの}，\ref{enu:真偽を問えないもの}，\ref{enu:真偽が定まらないもの}についてはのちほど説明するとして，ここでは\ref{enu:パラドックス}の例をあげよう．

\phantomsection
\begin{example}{自己言及のパラドックス}{自己言及のパラドックス}
次のような主張を考える：
\begin{quote}
この文は偽である．
\end{quote}
この文が真であるとすれば，その主張のとおり偽となってしまい，偽であるとすれば，主張どおりであるから真となってしまう．この文は\emph{嘘つきのパラドックス}として知られる自己言及のパラドックスであり，通常の二値論理の枠組みでは，そのまま命題として扱うことができない．
\end{example}

ここまで，二値論理では取り扱わない主張について述べてきたが，ここからは二値論理で取り扱う主張に戻ろう．

\phantomsection
\begin{example}{命題}{命題}
次に示す文は真の命題である：
\begin{enumerate}[(1)]
\item 「$1+1=2$」
\item 「$23 > 17$」
\item 「円周率は100未満である」
\item 「1日は24時間である」
\end{enumerate}
また，次に示す文は偽の命題である：
\begin{enumerate}[(a)]
\item 「$1+1=46$」
\item 「$23>26$」
\item 「円周率は3未満である」
\item 「1日は30時間である」
\end{enumerate}
\end{example}

\phantomsection
\begin{example}{命題でない文}{命題でない文}
次に示す文は命題ではない：
\begin{enumerate}[(A)]
\item 「$1+1$」（なにも主張しておらず，真か偽か判定できない）
\item 「この部屋は暑い」（「暑い」の基準が明確でなく，真か偽かが定まらない）
\item 「$10000$は大きい」（「大きい」の基準が明確でなく，真か偽かが定まらない）
\item 「この料理はおいしい」（「おいしい」という評価は人によって異なり，真か偽かが定まらない）
\item 「早く起きなさい」（命令文であって，真か偽か判定できない）
\item 「$x^2 > 4$」（$x$の値が定まっていないため，真か偽かが定まらない）\label{enu:命題でない文4}
\end{enumerate}
\end{example}

\phantomsection
\subsection{定理}

まず，前項で保留にしていた「狭い意味での命題」を確認しておこう．
数学では，証明によって正しいことが確かめられた主張に「命題」「定理」「補題」「系」といった見出し（ラベル）をつけて述べる．
「狭い意味での命題」とは，このラベルとしての「命題」のことであり，\textbf{Prop}と略記される．
前項でみた「真偽が定まる文」という広い意味と異なり，こちらは「証明された（あるいは証明すべき）正しい数学の主張」を指す言葉である\footnote{したがって，狭い意味での命題は，広い意味では「真の命題」である．偽である主張に「命題」というラベルがつけられることはない．}．
両者の違いを\tabref{命題の比較}に整理した．

\begin{table}[htbp]
  \centering
  \begin{tabular}{>{\centering\arraybackslash}m{2.2cm} L{4.9cm} L{4.9cm}}
    \toprule
      & \textbf{広い意味} & \textbf{狭い意味} \\
    \midrule
    なにを指すか & 真偽が定まっている文 & 証明された数学の主張につけるラベル \\
    真偽 & 真の場合も偽の場合もある & つねに真 \\
    対になる語 & 「命題でない文」 & 定理・補題・系 \\
    主な登場場面 & 論理（真理値表など） & 教科書の見出し \\
    例 & 「$1+1=2$」（真），「$1+1=46$」（偽） & \prref{逆行列の一意性} \\
    \bottomrule
  \end{tabular}
  \caption{広い意味・狭い意味の命題の比較}
  \label{tab:命題の比較}
\end{table}

そのうえで，\emph{定理}\index{ていり@定理}（\emph{theorem}）とは，正しいと分かっている数学の主張（狭い意味での命題）の中でもとりわけ重要なものを指す．\textbf{Thm}と略記される．


定理とは，定義や公理，すでに得られた結果から論理的に証明された数学的主張であり，議論を進めるうえでの基礎となる．定理は\partref{証明}で扱う「証明」によって裏付けられているため，採用した定義・公理のもとで正しいことが保証されている．

定理が重要な主張であるがゆえに，「余弦定理」・「加法定理」・「ハイネ・ボレルの被覆定理」など，固有の名前が与えられているものも存在する．
その固有の名前は，定理の主張の詳細，もしくは発見者の名前にちなんで付けられることが多い．たとえば，「三平方の定理\footnote{「ピタゴラスの定理」とも呼ばれる．}」や「コーシー・シュワルツの不等式」である．

\phantomsection
\begin{example}{三平方の定理}{三平方の定理}
  三平方の定理の主張は以下のようになる：
  \begin{quotation}
    直角三角形の斜辺の長さを$c$とし，他の二辺の長さを$a$，$b$としたとき，
    \[
      a^2+b^2=c^2
    \]
    が成り立つ．
  \end{quotation}
\end{example}

\phantomsection
\begin{example}{コーシー・シュワルツの不等式}{コーシー・シュワルツの不等式}
  コーシー・シュワルツの不等式は，内積空間における基本的な不等式であり，以下のように表される：
  \begin{quotation}
    任意の内積空間において，ベクトル$u$と$v$に対して，
    \[
      \abs{\langle u, v \rangle} \leqq \norm{u} \cdot \norm{v}
    \]
    が成り立つ．
  \end{quotation}
  この不等式は解析学や線型代数学など，多くの分野で重要な役割を果たす．
\end{example}

\phantomsection
\begin{example}{微分積分学の基本定理}{微分積分学の基本定理}
  この定理は微分と積分の関係を明らかにするものであり，以下のように述べられる：
  \begin{quotation}
    $f$を区間$[a, b]$上で連続な実数値関数とする．このとき，
    \[
      F(x) = \int_a^x f(t) \, dt
    \]
    と定義すると，$F$は$[a, b]$上で連続，$(a, b)$上で微分可能であり，
    \[
      F'(x) = f(x) \qquad (a < x < b)
    \]
    が成り立つ．
  \end{quotation}
  この定理により，積分と微分が互いに逆操作であることが示される．
\end{example}
\phantomsection
\subsection{補題・系}

\emph{補題}\index{ほだい@補題}（\emph{lemma}）とは，定理や命題を証明するために用いられる補助的な命題のことである．\textbf{Lem}と略記される．

一方，\emph{系}\index{けい@系}（\emph{corollary}）とは，すでに証明された定理や命題から比較的容易に導かれる結果のことである．\textbf{Cor}と略記される．

補題や系を説明するためには，関連するいくつかの命題や定理が必要である．以下に具体的な例を示す．


\begin{example}{（補題）ロルの定理}{（補題）ロルの定理}
  実数値関数$f$ が閉区間 $[a,b]$ で連続，開区間 $(a,b)$ で微分可能，かつ $f(a)=f(b)$ とする．
  このとき，
  \[
  f'(c)=0
  \]
  を満たす点 $c \in (a,b)$ が少なくとも1つ存在する．
\end{example}

\begin{tleftbar}
\begin{proof}
最大値・最小値の定理\footnote{ここでは，最大値・最小値の定理の成立を認めるものとする．}により，$f$は$[a,b]$ で最大値 $M$ と最小値 $m$ をとる．

\begin{enumerate}[(i)]
  \item \textbf{$f$ が定数の場合}：
  任意の $x \in [a,b]$ で $f'(x)=0$ なので，任意の $c \in(a,b)$ で
  $f'(c)=0$ となる．

  \item \textbf{$f$ が定数でない場合}：
  端点では $f(a)=f(b)$ だから，
  $M>f(a)$ または $m<f(a)$ が成り立つ．
  したがって，$f$ の最大値または最小値は内部点
  $c \in(a,b)$ にて達成される．
\end{enumerate}

$c$ で最大をとる場合を示す（最小のときは不等号が逆になる）．
$h>0$ とすると， $f(c+h) \leqq f(c)$だから
\[
\frac{f(c+h)-f(c)}{h} \leqq 0 ,\quad \therefore ~ f'_+(c) \coloneqq \lim_{h \to 0+} \frac{f(c+h)-f(c)}{h}\leqq 0
\]
である．
一方 $h<0$ では同様に
\[
\frac{f(c+h)-f(c)}{h}\geqq 0 , \quad \therefore~  f'_-(c)\coloneqq \lim_{h\to0-}\frac{f(c+h)-f(c)}{h}\geqq 0
\]
である．
$f$の点$c$における微分可能性により，$f'_+(c)=f'_-(c)=f'(c)$であるから，
\[
0\leqq f'(c)\leqq 0 , \quad \therefore ~f'(c)=0
\]
である．
このことから，直ちに主張が従う．
\end{proof}
\end{tleftbar}


\begin{example}{（定理）平均値の定理}{（定理）平均値の定理}
  実数値関数$f$ が $[a,b]$ で連続，$(a,b)$ で微分可能なら，
  ある $c \in(a,b)$ が存在して
  \[
    f'(c)=\frac{f(b)-f(a)}{b-a}
  \]
  が成り立つ．
\end{example}

\begin{tleftbar}
\begin{proof}
  関数$g$を
  \[
    g(x)=f(x)-\frac{f(b)-f(a)}{b-a}(x-a)
  \]
  と定義すると， $g(a)=g(b)$であり，なおかつ$g$ は $[a,b]$ で連続，$(a,b)$ で微分可能なので，
  ロルの定理より $(a,b)$ に $g'(c)=0$ を満たす $c$ がある．
  すると
  \[
    0=g'(c)=f'(c)-\frac{f(b)-f(a)}{b-a}
  \]
  が成り立つ．この式から，直ちに主張が従う．
\end{proof}
\end{tleftbar}

\begin{example}{（系1）$f'(x)=0$ なら $f$ は定数}{（系1）導関数ゼロなら定数}
  実数値関数$f$ が区間 $I$ で微分可能で $f'(x)=0$ （$x \in I$）とする．
  このとき $f$ は $I$ 上で一定である．
\end{example}

\begin{tleftbar}
\begin{proof}
  任意の $x<y$ に平均値の定理を適用すると，ある $c\in(x,y)$ があって
  \[
     \frac{f(y)-f(x)}{y-x}=f'(c)=0
  \]
  が成り立つ．よってこのとき$f(y)=f(x)$である．いま$x$，$y$ は任意なので $f$ は$I$上で一定である．
\end{proof}
\end{tleftbar}

\begin{example}{（系2）導関数の符号と単調性}{（系2）単調性の系}
  実数値関数$f$ が区間 $I$ で微分可能とする．
  \begin{enumerate}
    \item $f'(x)\geqq 0$（$x\in I$） なら $f$ は $I$ で単調増加関数である．
    \item $f'(x)\leqq 0$（$x\in I$） なら $f$ は $I$ で単調減少関数である．
  \end{enumerate}
\end{example}

\begin{tleftbar}
\begin{proof}
  (1) を示す．$x<y$ をとると平均値の定理により
  \[
     \frac{f(y)-f(x)}{y-x}=f'(c) \quad (c\in(x,y))
  \]
  となる．右辺は $0$以上，かつ $y-x>0$ だから $f(y)-f(x)\geqq 0$である．このことから直ちに主張が従う．
  (2) も同様である．
\end{proof}
\end{tleftbar}

ここまで補題・定理・系という用語を説明してきたが，最後にひとつ注意しておこう．
実は，これらのラベルの使い分けに数学的に厳密な基準はない．
どの主張を「定理」とよび，どの主張を「補題」とよぶかは，
\emph{著者がその文章の中で，その主張をどう位置づけたか}を表しているにすぎない．
 
実際，本節では平均値の定理を証明するための補助として「ロルの定理」を紹介したが，
世間では「ロルの補題」ではなく「ロルの定理」の名で通っている．
逆に，「補題」の名を冠しながら，多くの定理よりも重要な主張もある\footnote{たとえば，集合論の「ツォルンの補題」は，それぞれの分野で中心的な役割を果たす主張であるが，歴史的な経緯から「補題」とよばれ続けている．}．
また，ある教科書で「定理」とされている主張が，別の教科書では「命題」とされていることも珍しくない．
 
したがって，「補題だから重要でない」「命題だから定理より格下だ」と機械的に判断するのではなく，
その文章の中でその主張が果たしている役割に注目して読むことが大切である．

\begin{comment}
\phantomsection
\begin{example}{（補題）ユークリッドの補題}{（補題）ユークリッドの補題}
  ユークリッドの補題は，整数論における基本的な結果であり，素因数分解の一意性を証明する際に重要な役割を果たす．その主張は以下の通りである：

  \begin{quote}
    素数$p$が整数$a$と$b$の積$ab$を割り切るならば，$p$は$a$と$b$の少なくとも一方を割り切る．
  \end{quote}

\end{example}

\begin{tleftbar}
  \begin{proof}
    素数$p$が$ab$を割り切るとする．もし$p$が$a$を割り切らないならば，$\gcd(p, a) = 1$である．
    このとき，ベズーの等式より，ある整数$s$と$t$が存在して，$sp + ta = 1$が成り立つ．
    両辺に$b$を掛けると，$spb + tab = b$となる．
    左辺の$spb$は$p$で割り切れるが，$tab$は$p$で割り切れるため，
    右辺の$b$も$p$で割り切れる．したがって，$p$は$b$を割り切る．
  \end{proof}
\end{tleftbar}

この補題を用いて，素因数分解の一意性（算術の基本定理）を証明することができる．

\phantomsection
\begin{example}{自然数の約数の個数}{自然数の約数の個数}
  自然数$n$の正の約数の個数$d(n)$は，$n$の素因数分解に基づいて以下の式で与えられる：

  \begin{equation}
    d(n) = (e_1 + 1)(e_2 + 1) \dotsm (e_k + 1)
  \end{equation}

  ただし，$n$を素因数分解して$n = p_1^{e_1} p_2^{e_2} \dotsm p_k^{e_k}$と表した．
\end{example}

\begin{tleftbar}
  \begin{proof}
    各素因数$p_i$について，指数$e_i$は$0$から$e_i$までの値を取り得る．
    したがって，各$p_i$に対する約数の取り得る指数の個数は$e_i + 1$個である．
    すべての素因数について独立に指数を選ぶことで，
    $n$のすべての約数を生成できるため，約数の総数は各$(e_i + 1)$の積となる．
  \end{proof}
\end{tleftbar}

この命題は，素因数分解の一意性から直接的に導かれる結果である．

そしてこの例から，次の系が導かれる：

\phantomsection
\begin{example}{（系）約数の個数が奇数となる条件}{（系）約数の個数が奇数となる条件}
  自然数$n$の約数の個数$d(n)$が奇数となるための必要十分条件は，$n$が完全平方数であることである．
\end{example}

\begin{tleftbar}
  \begin{proof}
    約数の個数$d(n) = (e_1 + 1)(e_2 + 1) \dotsm (e_k + 1)$である．
    各$(e_i + 1)$が奇数となるためには，$e_i$が偶数でなければならない．
    つまり，すべての素因数の指数$e_i$が偶数であるとき，$n$は各素因数の偶数乗の積であり，
    これは$n$が完全平方数であることと同値である．
    逆に，$n$が完全平方数であれば，各$e_i$は偶数であり，したがって$d(n)$は奇数となる．
  \end{proof}
\end{tleftbar}
\end{comment}


\phantomsection
\section{数学でよく使われる言い回し}

数学の文章では，日常語と同じ単語が，日常とは少しずれた意味や役割で使われることが多い．
「かつ・または」のような論理を表す言葉から，「嬉しい」「天下り」のような数学者特有の口ぐせまで，
数学書を読み始めた学生が引っかかりやすい言い回しを，具体例とともにひとつずつ確認していこう．

\phantomsection
\subsection{かつ・または}

数学において，「かつ」は日常とほとんど同じ意味で使われるが，「または」の使い方は日常とは異なる．以下に例を挙げよう．

\phantomsection
\begin{example}{「または」の使用例}{「または」の使用例}
  \begin{enumerate}[(A)]
    \item ランチメニューの主食として，米またはパンがついてくる\footnote{一般的には，この文を「米とパンの両方が食べられる」とは解釈しないであろう．}．\label{enu:米またはパン}
    \item 「運転免許を持っていない人」または「18歳未満の人」はレンタカーを借りることができない． \label{enu:運転免許または18歳未満}
  \end{enumerate}
  \ref{enu:米またはパン}は「どちらか片方のみ」の意味で「または」を使い，\ref{enu:運転免許または18歳未満}は「いずれかが」の意味で「または」を用いている．
\end{example}

数学では，「または」は「いずれかが」の意味で使われ，「どちらか片方のみ」の意味で使われることはない．

たとえば，$A$，$B$を集合とするとき，
\[
  x \in A \cup B
\]
は「$x$は$A$の元であるか，$B$の元であるか，あるいは両方である」という意味である\footnote{和集合$\cup$の定義は\partref{集合}で改めて述べる．}．
つまり，「$ x \in A$であり，$x \notin  B$である」あるいは「$ x \notin A$であり，$x \in B$である」といった状況のときにも，$x \in A \cup B$とかく．
%ここまで

\subsection{任意の・すべての}

数学において，「任意の\index{にんいの@任意の}」と「すべての\index{すべての@すべての}」はニュアンスが異なる．
英語では，「任意の」は``for any''，「すべての」は``for all''に相当する．
``for any''のニュアンスは「複数の対象があり，その中のどのひとつをとっても」という意味であり，
``for all''は「対象全体を見て，すべてのものが」という意味である．

たとえば，次の文はいずれも自然な日本語である．
\begin{itemize}
  \item 任意の実数$x$について，$x^2\ge0$である．
  \item すべての実数$x$について，$x^2\ge0$である．
\end{itemize}

一方，証明では，
\begin{quote}
任意に$x\in A$をとる．
\end{quote}
のような表現をよく用いる．
これは「集合$A$の中からどの元を選んでもよい」という意味であり，
証明の出発点として自然な表現である．

これに対して，
\begin{quote}
すべての$x\in A$をとる．
\end{quote}
という表現は，日本語として不自然である．
「すべての元」を同時にひとつずつ選ぶことはできないからである．

このように，「任意の」と「すべての」は論理的には同じ意味で用いられることが多いが，
日本語としては役割が異なるため，状況に応じて使い分ける必要がある．
例としてまとめよう：

\begin{example}{「任意の」と「すべての」の使用例}{「任意の」と「すべての」の使用例}
  \begin{enumerate}[(A)]
    \item 任意の実数$x$について，$x^2\ge0$である．
    \item すべての実数$x$について，$x^2\ge0$である．
    \item 任意に$x\in A$をとる．
  \end{enumerate}
\end{example}

\subsection{とる\index{とる}・おく\index{おく}・固定する\index{こていする@固定する}}

数学書の証明を読み始めると，
\begin{quote}
  任意に$x \in A$をとる．
\end{quote}
という表現に必ず出会う．初めて読むと「$x$をどこから取ってくるの？」と，かなり奇妙な日本語に感じられるかもしれない．
これは「集合$A$の元をひとつ選び，それを$x$と名づけて議論の対象にする」という意味であり，「選ぶ」に近い言葉である．
「任意に」がついている場合は，前項で述べたとおり「どの元を選んでも以下の議論が通用する」というニュアンスが込められている．
また，「十分大きな$n$をとる」のように，条件を満たす対象を選ぶ場面でも使われる．

これとよく似た言葉に「おく」がある．
\begin{quote}
  $s=x+y$とおく．
\end{quote}
のように，「おく」は既にある量や式に対して\emph{新しい記号を導入する}（名前をつける）際に使う言葉である．
「とる」が対象を選ぶ操作であるのに対し，「おく」は記号の定義・命名の操作である点が異なる．

もうひとつ，「固定する」という言葉もある．
\begin{quote}
  $x \in A$をひとつ固定する．
\end{quote}
これは「$x$をひとつ選び，\emph{以後の議論ではそれを動かさない}」という宣言である．
複数の文字が登場する議論で，「$x$は止めておいて，$y$だけを動かして考える」といった場面で威力を発揮する．
たとえば「$x$を固定するごとに，$y$についての関数$f(x,y)$を考える」のように使われる．

\begin{example}{「とる」「おく」「固定する」の使用例}{「とる」「おく」「固定する」の使用例}
  \begin{enumerate}[(A)]
    \item 任意に$x \in A$をとる．（$A$の元をひとつ選ぶ．どの元でも以下の議論が通用する）
    \item 十分大きな$n \in \mathbb{N}$をとる．（条件を満たす$n$を選ぶ）
    \item $t = \tan \frac{x}{2}$とおく．（新しい記号$t$を導入する）
    \item $\varepsilon > 0$をひとつ固定する．（以後の議論で$\varepsilon$は動かさない）
  \end{enumerate}
\end{example}

この3語は，証明の冒頭で文字を導入するときの基本装備である．
「とる」は選ぶ，「おく」は名づける，「固定する」は動かさないと宣言する，と整理して覚えておくとよい．

\subsection{成り立つ\index{なりたつ@成り立つ}}

数学において「成り立つ」は，等式・不等式・命題などが\emph{真である}ことを表す言葉である．
日常で「商売が成り立つ」「会話が成り立つ」というときの「うまく機能する」という意味とは少し違い，
数学では「その主張が正しい」ということを端的に表す．

\begin{example}{「成り立つ」の使用例}{「成り立つ」の使用例}
  \begin{enumerate}[(A)]
    \item 任意の実数$x$に対して，不等式$x^2 \geqq 0$が成り立つ．
    \item $n$が偶数のとき，$n^2$も偶数である，という命題が成り立つ．
  \end{enumerate}
\end{example}

主語になるのは「等式」「不等式」「命題」といった主張そのものであることに注意しよう．
本書でもすでに「$a^2+b^2=c^2$が成り立つ」のような形で繰り返し使ってきたが，
すべて「この式は真である」という意味である．

\subsection{一般に\index{いっぱんに@一般に}}

「成り立つ」とセットでよく使われる言葉に「一般に」がある．
これは，肯定文と否定文とで読み方の注意点が異なる，少し曲者の言い回しである．

まず，肯定文とともに使われる場合を見よう．
\begin{quote}
  一般に，実数$x$に対して$x^2 \geqq 0$が成り立つ．
\end{quote}
このときの「一般に」は「\emph{すべての場合に}」という意味であり，
「任意の実数$x$に対して成り立つ」という全称の主張と同じことを述べている．

要注意なのは，否定文とともに使われる場合である．
\begin{quote}
  $\sqrt{a^2} = a$は一般に成り立たない．
\end{quote}
これは「\emph{つねに成り立つとは限らない}」，すなわち「成り立たない場合が少なくともひとつ存在する」という意味である．
決して「すべての場合に成り立たない」という意味ではない．
実際，$\sqrt{a^2}=a$は$a \geqq 0$のときには成り立ち，$a < 0$のときに成り立たない
（\exref{「絶対値」の定義}で見たとおり，正しくは$\sqrt{a^2}=\abs{a}$である）．
このように，成り立つ場合がいくらあっても，例外がひとつでもあれば「一般に成り立たない」というのである．
「一般には成り立たない」と，「は」を入れて書かれることもあるが，意味は同じである．

\begin{example}{「一般に」の使用例}{「一般に」の使用例}
  \begin{enumerate}[(A)]
    \item 一般に，$(a+b)^2 = a^2 + 2ab + b^2$が成り立つ．（任意の実数$a$，$b$で成り立つ）
    \item 正方行列$A$，$B$に対して，$AB=BA$は一般に成り立たない．（成り立つ組もあるが，成り立たない組が存在する）\label{enu:積の非可換}
  \end{enumerate}
\end{example}

「一般に成り立たない」ことを証明するには，成り立たない例，すなわち\emph{反例}\index{はんれい@反例}（\emph{counterexample}）をひとつ挙げれば十分である．
たとえば\ref{enu:積の非可換}については，
\[
  A=\begin{pmatrix} 1 & 1 \\ 0 & 1 \end{pmatrix},\quad
  B=\begin{pmatrix} 1 & 0 \\ 1 & 1 \end{pmatrix}
\]
とすれば，$AB \ne BA$となることが計算により確かめられる\footnote{一方，$B=E$とすれば$AB=BA$である．「成り立つ例がある」ことと「一般に成り立つ」ことは別物である．}．
全称の主張（一般に成り立つ）を示すにはすべての場合を扱う議論が必要である一方，
その否定（一般に成り立たない）は反例ひとつで示せる．この非対称性は，数学の議論の基本として覚えておきたい．
なお，のちに述べる「一般性を失わず」の「一般性」も，この「すべての場合」という意味での「一般」を指している．

\subsection{示す\index{しめす@示す}}

数学において「示す」は「証明する」とほぼ同じ意味で使われる．
演習問題などで頻出の「〜を示せ」は「〜を証明せよ」という意味である．

また，証明の中では
\begin{quote}
  (1)を示す．
\end{quote}
のように，\emph{これから証明する目標を宣言する}ためにも使われる．
長い証明では「まず$\subseteq$を示す」「次に$\supseteq$を示す」のように，
証明全体をいくつかの目標に分割する際の目印となる．
証明を読むときは，「示す」と宣言された目標がどこで達成されたかを意識すると，議論の構造が見えやすくなる．

\begin{example}{「示す」の使用例}{「示す」の使用例}
  \begin{enumerate}[(A)]
    \item $\sqrt{2}$が無理数であることを示せ．（＝証明せよ）
    \item まず，数列$(x_n)_{n \in \mathbb{N}}$が上に有界であることを示す．（＝これから証明する）
  \end{enumerate}
\end{example}

\subsection{仮定する\index{かていする@仮定する}}

数学において「仮定する」は，ある主張を\emph{一時的に正しいものとみなして}議論を進めることを表す．
真であると断定しているわけではないことに注意しよう．

「仮定する」が活躍する典型的な場面はふたつある．
ひとつは，数学的帰納法のように「$n=k$のとき成り立つと仮定する」と述べて，そこから$n=k+1$の場合を導く場面である．
もうひとつは背理法であり，「$\sqrt{2}$が有理数であると仮定する」のように，
\emph{最終的には否定したい主張}をあえて一時的に認め，矛盾を導くために使われる．

\begin{example}{「仮定する」の使用例}{「仮定する」の使用例}
  \begin{enumerate}[(A)]
    \item $n=k$のとき主張が成り立つと仮定する．（数学的帰納法）
    \item $\sqrt{2}$が有理数であると仮定する．（背理法：のちに矛盾を導く）
  \end{enumerate}
\end{example}

いずれの場合も，「仮定」はその議論の中だけで有効な約束であり，議論が終われば効力を失う．
「仮定する」と書かれていたら，「この仮定はどの範囲で使われ，どこで役目を終えるのか」を意識して読むとよい．

\subsection{十分大きい\index{じゅうぶんおおきい@十分大きい}}

日常語の「十分大きい」と，数学における「十分大きい」には，かなりの差がある．
数学書で
\begin{quote}
  十分大きな$n$に対して，$a_n < \varepsilon$が成り立つ．
\end{quote}
と書かれていたら，これは「$n$がなんとなく大きければ」という曖昧な意味ではない．正確には
\begin{quote}
  ある$N \in \mathbb{N}$が存在して，$n \geqq N$を満たす\emph{すべての}$n$について$a_n < \varepsilon$が成り立つ．
\end{quote}
という意味である．
つまり「十分大きい」は，「ある値（しきい値）以上ならすべて」という，存在と全称を組み合わせた厳密な主張の省略表現なのである．

\begin{example}{「十分大きい」の使用例}{「十分大きい」の使用例}
  \begin{enumerate}[(A)]
    \item 十分大きな$n$に対して，$\dfrac{1}{n} < 0.01$が成り立つ．
    （実際，$N=101$とすれば，$n \geqq N$なるすべての$n$で成り立つ）
    \item 十分大きな$x$に対して，$x^2 > 100x$が成り立つ．
    （$x > 100$なるすべての$x$で成り立つ）
  \end{enumerate}
\end{example}

この言い回しは，のちに学ぶ数列の極限（$\varepsilon$-$N$論法）に直結する．
「十分大きい」を見たら「しきい値がどこかに存在する」と読み替える習慣をつけておくと，極限の定義がぐっと理解しやすくなる．

\subsection{簡単のため\index{かんたんのため@簡単のため}}

「簡単にするために」のほうがしっくりくる方もおられるかもしれないが，これは``For simplicity''の訳で，
その後の議論を簡単にするために「議論することが楽な仮定」を設ける際に使う言葉である．

\begin{example}{「簡単のため」の使用例}{「簡単のため」の使用例}
  \begin{enumerate}[(A)]
    \item 簡単のため，平行移動して頂点が原点にある形として，$f(x)=ax^2$を考える．
  \end{enumerate}
\end{example}

\subsection{一般性を失わず\index{いっぱんせいをうしなわず@一般性を失わず}}

数学書には
\begin{quote}
  一般性を失わず，$a \leqq b$としてよい．
\end{quote}
という表現がしばしば登場する．初めて読むと「勝手に$a \leqq b$と仮定していいの？」と戸惑うかもしれない．
これは，
\begin{quote}
  残りの場合（ここでは$a > b$の場合）も，文字を入れ替えるなどして同じ議論で処理できるので，一方の場合だけ考えても問題がない
\end{quote}
という意味である．英語では``without loss of generality''といい，\textbf{WLOG}と略記されることもある．

\begin{example}{「一般性を失わず」の使用例}{「一般性を失わず」の使用例}
  実数$a$，$b$に対して，$\max\{a,b\} + \min\{a,b\} = a+b$が成り立つことを示したい．
  このとき，
  \begin{quote}
    一般性を失わず，$a \leqq b$としてよい．すると$\max\{a,b\}=b$，$\min\{a,b\}=a$なので，
    \[
      \max\{a,b\} + \min\{a,b\} = b + a = a + b
    \]
    が成り立つ．
  \end{quote}
  と議論できる．$a > b$の場合は，$a$と$b$の役割を入れ替えれば同じ議論がそのまま通用するので，
  改めて書く必要がないのである（示したい等式が$a$と$b$について対称であることが効いている）．
\end{example}

「一般性を失わず」は，数学特有の\emph{省略の仕方}のひとつである．
ただし，本当に残りの場合が同じ議論で処理できるのかは，読者・書き手の双方が確認すべきことである．
対称性がないのに「一般性を失わず」と書いてしまうと，証明に穴があくので注意しよう．

\subsection{同様に\index{どうように@同様に}}

「同様に」も，数学特有の省略の表現である．
\emph{本当に同じ論法がそのまま使える}部分について，繰り返しを避けるために「同様に」と述べて詳細を省略する．

実際，本書でも「導関数の符号と単調性」の系の証明において，(1)を証明したあとに
\begin{quote}
  (2)も同様である．
\end{quote}
と述べて証明を終えた．これは「(1)の議論の不等号の向きを逆にすれば，そのまま(2)の証明になる」という意味である．

\begin{example}{「同様に」の使用例}{「同様に」の使用例}
  \begin{enumerate}[(A)]
    \item $c$で最大をとる場合を示す（最小のときは不等号が逆になるだけで，同様である）．
    \item $x \geqq 0$の場合は上で示した．$x<0$の場合も同様に示される．
  \end{enumerate}
\end{example}

「同様に」と書かれた箇所は，著者が「読者が自力で再現できる」と判断して省略した部分である．
数学書を読むときは，「同様に」の中身を実際に自分の手でなぞって確認してみると，理解の確認になり，力もつく．

\subsection{従う\index{したがう@従う}}

数学において「従う」は日常とはまた違った意味で使われる．日常だと「王に従う」や「上司に従う」など，「命令を受けてそれに従う」という意味で使われることが多いが，
数学では「$A$という事柄から，すぐ$B$という事柄がわかる」という場合に「$A$から$B$が従う」という．

\begin{example}{「従う」の使用例}{「従う」の使用例}
  \begin{enumerate}[(A)]
    \item $n$が偶数であることから，$n^2$が偶数であることが従う．\label{enu:偶数ならば偶数}
    \item $G$が群であることから，$G \ne \varnothing$が従う\footnote{群であれば単位元が存在することは群の定義からわかる．そのため群は空集合でない．}．\label{enu:群の例}
  \end{enumerate}
\end{example}

\subsection{明らか\index{あきらか@明らか}・直ちに\index{ただちに@直ちに}}

数学書には「明らかである」「明らかに〜が成り立つ」という表現が頻出する．
文字どおり受け取ると「証明不要なほど自明」という意味に思えるが，実際には
\begin{quote}
  ここでは詳細を省略できる程度の主張である（本筋から外れるので，証明は読者に委ねる）
\end{quote}
という意味で使われることが多い．
つまり「明らか」は主張の真偽についての言葉というより，\emph{著者の省略の宣言}なのである．
そのため，読者にとってはまったく明らかでないこともしばしばある\footnote{「明らか」と書かれた箇所でつまずくのは恥ずかしいことではない．むしろ，行間を自力で埋める良い練習問題だと考えるとよい．}．

一方，「直ちに」は，既知の事実から\emph{ごく短い推論で}結論が得られることを表す．
実際，本書でもロルの定理の証明などで
\begin{quote}
  このことから，直ちに主張が従う．
\end{quote}
という表現を用いた．これは「得られた式から，ほんの一，二歩の変形・言い換えで結論に到達する」という意味である．

\begin{example}{「明らか」「直ちに」の使用例}{「明らか」「直ちに」の使用例}
  \begin{enumerate}[(A)]
    \item $x > 0$は明らかである．（詳細は省略するが，各自で確認できる程度の主張である）
    \item $f'(c)=0$から，直ちに主張が従う．（短い推論で結論が得られる）
  \end{enumerate}
\end{example}

「明らか」「直ちに」に出会ったら，鵜呑みにせず，省略された推論を自分の手で復元してみることを勧める．
それが「本当に明らか」であれば数行で終わるし，そうでなければ，そこが理解の急所である．

\subsection{〜によらない\index{によらない@〜によらない}}

数学において「〜によらない」は，\emph{選び方・表し方・添字などを変えても結果が変わらない}ことを表す言葉である．

\begin{example}{「〜によらない」の使用例}{「〜によらない」の使用例}
  \begin{enumerate}[(A)]
    \item この関数の値は，有理数の表し方によらず定まる．
    \item 三角形の内角の和は，三角形の選び方によらず$180^\circ$である．
    \item 極限値$\lim_{n \to \infty} a_n$は，最初の有限項の値によらない．
  \end{enumerate}
\end{example}

この言葉は，本章の「定義」の項で述べたwell-definedの概念と深く関係している．
「$g\left(\frac{a}{b}\right)=\frac{a}{2b}$という定義は\emph{分数の表し方によらず}値を定める」からこそ，$g$はwell-definedなのであった．
数学では，なにかを定義したときに「その定義が途中の選択によらないこと」の確認が要求される場面が多い．
「〜によらない」という表現を見たら，「なにを取り替えても結果が同じなのか」を意識して読むとよい．

\subsection{〜に関して\index{にかんして@〜に関して}}

数学において「〜に関して」は，\emph{どの文字（変数）に注目して操作・議論を行うのか}を指定する言葉である．

\begin{example}{「〜に関して」の使用例}{「〜に関して」の使用例}
  \begin{enumerate}[(A)]
    \item 関数$f(x,y)=x^2 y$を$x$に関して微分すると，$2xy$である．
    （$y$は定数とみなし，$x$を変数として微分する）
    \item $f(x)=x^2$は，$x \geqq 0$の範囲で$x$に関して単調増加である．
    \item 等式$x+y=y+x$は，$x$と$y$に関して対称である．
  \end{enumerate}
\end{example}

複数の文字が登場する式では，「どの文字を動かし，どの文字を止めているのか」が曖昧だと議論が混乱する．
「$x$に関して」という指定は，先に述べた「固定する」と対になる表現であり，
「$y$は固定して，$x$に関して微分する」のように組み合わせて使われることも多い．

\subsection{嬉しい\index{うれしい@嬉しい}}

日常では，「嬉しい」という言葉は「喜びを感じる」という意味で使われることが多いが，数学の文脈では「都合がよい」，「議論が楽になる」という意味でたびたび使われる．

\begin{example}{「嬉しい」の使用例}{「嬉しい」の使用例}
  \[
    \int_{\alpha}^{\beta} (x-\alpha)(x-\beta)\, dx = -\frac{(\beta-\alpha)^3}{6}
  \]
  であることを示したい．このときに
  \[
    \int_{\alpha}^{\beta} (x-\alpha)(x-\beta)\, dx = \int_{\alpha}^{\beta}  (x-\alpha)\{(x-\alpha) - (\beta-\alpha)\} \, dx
  \]
  と変形してなにが嬉しいかというと，
  \begin{align*}
    \int_{\alpha}^{\beta}  (x-\alpha)\{(x-\alpha) - (\beta-\alpha)\} \, dx & = \int_{\alpha}^{\beta} \{  (x-\alpha)^2 -(\beta-\alpha)(x-\alpha) \} \, dx                        \\
                                                                           & = \Biggl [ \frac{(x-\alpha)^3}{3} - \frac{(\beta-\alpha)(x-\alpha)^2}{2} \Biggr ]_{\alpha}^{\beta} \\
                                                                           & = - \frac{1}{6} (\beta-\alpha)^3
  \end{align*}
  というふうに，簡単に計算できるからである．
\end{example}



\subsection{おさえる\index{おさえる}・評価\index{ひょうか@評価}する}

数学の文脈では「おさえる」という言葉は，ある数$K$を，別の数 $M$を用いて，$K<M$という形で表すことを指す．
より詳しく「上からおさえる」と表すこともある．

また，数学の文脈で「評価」という言葉は，ある数$M$を不等式で表し，その数がどの程度の大きさであるかを示すことを指す．
たとえば，$ 1 < M <4$と表すことは「$M$を$1$より大きく$4$より小さいと評価する」ということを指す．
さらに，上記の$M$を$ 2< M <3$と表すことができる場合は，後者の表現のほうが「評価が厳しい」ということができる．

\begin{example}{「おさえる」と「評価する」の使用例}{「おさえる」と「評価する」の使用例}
  \begin{enumerate}[(A)]
    \item 数列$(a_n)_{n \in \mathbb{N}}$を $ a_n <M$と上からおさえる．
    \item $e$の値を$2.7 < e < 2.8$と評価する．
  \end{enumerate}
\end{example}


\subsection{特徴づけ\index{とくちょうづけ@特徴づけ}}

数学の文脈では，「特徴づけ」（\emph{characterization}）という言葉がよく使われる．
これは，ある対象（あるいはあるクラスに属すること）と\emph{同値な条件を与える}ことを指す．
「$X$であることと，条件$P$を満たすことは同値である」という形の主張が得られたとき，「条件$P$は$X$を特徴づける」というのである．

\begin{example}{「特徴づけ」の使用例}{「特徴づけ」の使用例}
  \begin{enumerate}[(A)]
    \item 偶数は，「ある整数$k$を用いて$n=2k$と表せる整数$n$」として特徴づけられる．
    （$n$が偶数であることと，$n=2k$と表せることは同値である）
    \item 正則行列は，「行列式が$0$でない正方行列」として特徴づけられる\footnote{これは線型代数で学ぶ重要な定理である．定義（逆行列を持つこと）とは別の，判定しやすい同値条件が得られている点に価値がある．}．
    \item 二等辺三角形は，「ふたつの内角が等しい三角形」として特徴づけられる．
  \end{enumerate}
\end{example}

特徴づけの価値は，\emph{定義とは別の切り口から同じ対象を捉え直せる}ことにある．
たとえば偶数の定義を「$2$で割り切れる整数」としたとき，「$n=2k$と表せる」という特徴づけは，
偶数に関する等式の計算を進めるうえで扱いやすい形を与えてくれる．
定義・特徴づけのどちらを使って議論してもよいが，そのためには両者が同値であることの証明が必要である
（これは本章の冒頭で述べた，絶対値のふたつの定義の同値性の話と同じ構図である）．

\subsection{天下り\index{あまくだり@天下り}}

数学の文脈で，数式や概念を根拠を明示せずに導入することを「天下り」とよぶことがある．

\begin{example}{「天下り」の使用例}{「天下り」の使用例}
  \begin{enumerate}[(A)]
    \item $a$，$b$を$a<b$を満たす実数とするとき，$ a < c < b$となる$c$の存在を示すために，
    天下りではあるが，$c=(a+b)/2$と定める\label{enu:天下りの例}．
  \end{enumerate}
\end{example}

\ref{enu:天下りの例}においては，根拠を明示せずに$c=(a+b)/2$としたが，これは「$c$は$a$と$b$の相加平均である」という根拠に基づいて定めている．
% NOTE: この部は執筆中（今後，節を追記する可能性あり）